\par\noindent\SourceLeaf{49}%
\hspace*{\parindent}Das allgemeine Glied $r^{\text{tus}}$
\[
= \frac{1}{2^{2m-1}}\,{}^{2m}\mathfrak{B}^{(m+r)}\cos 2r\varphi\cos\varphi,
\]
\par\noindent ergibt, in Teile zerlegt,
\[
\frac{1}{2^{2m-1}}\,{}^{2m}\mathfrak{B}^{(m+r)} . \frac{1}{2}\cos(2r-1)\varphi
+ \frac{1}{2^{2m-1}}\,{}^{2m}\mathfrak{B}^{(m+r)} . \frac{1}{2}\cos(2r+1)\varphi.
\]

\par\noindent Hier treten also nur ungerade Vielfache auf, von denen je zwei benachbarte aus jedem Glied der gegebenen Reihe entstehen. Sucht man daher das Glied $(r-1)^{\text{tum}}$ der neuen Reihe, das offenbar mit $\cos(2r-1)\varphi$ multipliziert ist, so nimmt man aus dem Glied $(r-1)^{\text{to}}$ der entwickelten Reihe den zweiten Teil
\[
= \frac{1}{2^{2m}}\,{}^{2m}\mathfrak{B}^{(m+r-1)}\cos(2r-1)\varphi,
\]
und aus dem Glied $r^{\text{to}}$ den ersten Teil
\[
\frac{1}{2^{2m}}\,{}^{2m}\mathfrak{B}^{(m+r)}\cos(2r-1)\varphi
\]
und erhält durch Addition
\[
\frac{1}{2^{2m}}\,{}^{2m+1}\mathfrak{B}^{(m+r)}\cos(2r-1)\varphi.
\]

\par\noindent So nimmt für $\cos\varphi^{2m+1}$ das allgemeine Glied $(r-1)^{\text{tus}}$ die hier vorgeschlagene Form an. Oder man erhält, wenn man lieber für $r$ den Ausdruck $r+1$ einsetzt, das Glied
$r^{\text{tus}}$
\[
\frac{1}{2^{2m}}\,{}^{2m+1}\mathfrak{B}^{(m+1+r)}\cos(2r+1)\varphi.
\]

\par\noindent Es ist nun klar, dass dieses Gesetz, das für $\cos\varphi^{2m+1}$ das Glied $r^{\text{tum}}$ darstellt, mit dem für $\cos\varphi^{2m-1}$ angenommenen Gesetz für das Glied $r^{\text{to}}$
\[
= \frac{1}{2^{2m-2}}\,{}^{2m-1}\mathfrak{B}^{(m+r)}\cos(2r+1)\varphi,
\]
so übereinstimmt, dass sich das erstere unmittelbar ergibt, sobald man im letzteren $m$ durch $m+1$ ersetzt.

\par\noindent\hspace*{\parindent}Da dieses unbegrenzte Fortschreiten durch den Beweis gesichert ist, liegt auf der Hand, dass das Gesetz sowohl für gerade als auch für ungerade Potenzen von $\cos\varphi$ allgemein gilt.

5. Obwohl es nicht zu der vorliegenden Untersuchung gehört, sei hier bemerkt, dass sich aus dem Gesetz, das die Potenzen des Kosinus durch Vielfache des gegebenen Winkels darstellt, unmittelbar ein entsprechendes Gesetz für die Potenzen des Sinus ergibt. Dazu braucht man nur zu beachten, dass
$\cos(90^0-\varphi) = \sin\varphi$.
